%%%%%%%%%%  scan idx 448  |  folio 435  %%%%%%%%%%
Pour vérifier que la suite est exacte, on peut le faire localement sur \(T\), et utiliser les sections définies dans l'énoncé de 2.5. On a alors une base de \(N_{T/X}\) formée des images des \(f_i,g_j,h_k\), une base de \(i^*(N_{Y/X})\) formée des images des \(f_i,g_j\), et une base de \(j^*(N_{Z/X})\) formée des images des \(f_i,h_k\). Il est alors clair que \(i^*(N_{Y/X})\times j^*(N_{Z/X})\to N_{T/X}\) est surjectif, et que son noyau est formé des éléments ayant des coordonnées nulles sur les \((\overline g_j,0)\) et les \((0,\overline h_k)\), et des coordonnées égales pour tout \(i\) sur \((\overline f_i,0)\) et \((0,\overline f_i)\). Par suite, la suite est également exacte au milieu. Or les trois \(\cO_T\)-Modules intervenant sont localement libres de type fini; comme \(\underline{\Tor}_1^{\cO_X}(\cO_Y,\cO_Z)\) est de rang \(p\) d'après la démonstration du i), on voit immédiatement que la suite est exacte.

\subsection*{2.6}
On suppose maintenant que \(X\) est un schéma ayant un faisceau ample. Tout sous-schéma fermé \(Y\) de \(X\) a alors un faisceau ample. On sait que dans ce cas ``il existe des résolutions globales'', et que les \(K^\circ\) formés respectivement avec les faisceaux localement libres de type fini et les complexes parfaits sont identiques (II, 2.2.9). Si en outre l'immersion \(i:Y\to X\) est régulière, c'est un morphisme parfait, et on peut définir un homomorphisme additif \(i_*:K^\circ(Y)\longrightarrow K^\circ(X)\) (IV 2.12), en associant à la classe d'un complexe parfait \(L^\bullet\) la classe du complexe \(\mathbb Ri_*(L^\bullet)=i_*(L^\bullet)\) (\(i_*\) étant un foncteur exact).

\textbf{Proposition 2.7.} Soient \(X\) un schéma ayant un faisceau ample, \(i:Y\to X\) une immersion fermée régulière, \(N\) le faisceau conormal de \(Y\) dans \(X\). Alors, pour tout \(x\in K^\circ(Y)\), on a la relation:
\begin{equation}\notag
i^*(i_*(x))=x.\lambda_{-1}(N),
\end{equation}
où \(N\) est la classe de \(N\) dans \(K^\circ(Y)\), et \(\lambda_{-1}(N)=\sum_{i\ge0}(-1)^i\lambda^i(N)\).


